% SOURCES: docs/diagrams/*.png (already in report/figures/); docs/0.2.0v/architecture.md;
%          docs/0.2.0v/element-model.md (the class model); src/document/types.ts (ground
%          truth for the class diagram); docs/1.0.0v/supervisor-report.md (§2 diagram)
% GUIDANCE: The marking scheme asks for UML (class, sequence, use-case). Five
% architecture figures already exist and are wired in below. You still need to
% PRODUCE three formal UML diagrams — PlantUML skeletons are provided as comments
% so you can render them to figures/ and \includegraphics them.

\chapter{System Analysis}
\label{ch:analysis}

\section{Architectural Overview}
% DRAFT (review): the existing diagrams carry this section.
Draw-to-Web is organised into layers with a strictly one-directional data flow:
the user interface writes to the store through operations, the document layer is
the single source of truth, and the output layer reads the document to generate
and export. Figure~\ref{fig:architecture} shows the layered architecture and
Figure~\ref{fig:dataflow} the one-directional data flow.

\begin{figure}[H]
  \centering
  \includegraphics[width=0.92\linewidth]{01-architecture.png}
  \caption{Layered architecture: UI $\rightarrow$ Document Model $\rightarrow$ output pipeline.}
  \label{fig:architecture}
\end{figure}

\begin{figure}[H]
  \centering
  \includegraphics[width=0.92\linewidth]{02-dataflow.png}
  \caption{One-directional data flow. The canvas renders the document; the generator walks it.}
  \label{fig:dataflow}
\end{figure}

\begin{figure}[H]
  \centering
  \includegraphics[width=0.88\linewidth]{04-component-hierarchy.png}
  \caption{Renderer component hierarchy.}
  \label{fig:components}
\end{figure}

\begin{figure}[H]
  \centering
  \includegraphics[width=0.80\linewidth]{05-state-stores.png}
  \caption{State stores and their relationships.}
  \label{fig:stores}
\end{figure}

\section{Class Diagram (Document Model)}
% SOURCES: src/document/types.ts; docs/0.2.0v/element-model.md
% GUIDANCE: The Document Model is the class model. Generate a UML class diagram
% from src/document/types.ts and place it here. PlantUML starter (render to
% figures/class-diagram.png):
%
%   @startuml
%   class Document {
%     +version: string
%     +tree: ElementNode
%     +tokens: Tokens
%     +seo: SEOConfig
%     +runtime: RuntimeFlags
%     +variables: DocumentVariables
%   }
%   abstract class ElementNode {
%     +id: string
%     +semanticRole
%     +style: ResponsiveProperties
%     +states: StatesMap
%   }
%   class ContainerNode
%   class TextNode
%   class ImageNode
%   class ButtonNode
%   class LinkNode
%   class IconNode
%   class ListNode
%   class DividerNode
%   Document "1" *-- "1" ElementNode : tree
%   ElementNode <|-- ContainerNode
%   ElementNode <|-- TextNode
%   ElementNode <|-- ImageNode
%   ElementNode <|-- ButtonNode
%   ElementNode <|-- LinkNode
%   ElementNode <|-- IconNode
%   ElementNode <|-- ListNode
%   ElementNode <|-- DividerNode
%   ContainerNode "1" o-- "*" ElementNode : children
%   Document *-- Tokens
%   @enduml
%
\begin{figure}[H]
  \centering
  % TODO: \includegraphics[width=\linewidth]{class-diagram.png}
  \fbox{\parbox[c][5cm][c]{0.9\linewidth}{\centering\itshape
    TODO: class diagram of the Document Model, generated from
    \texttt{src/document/types.ts} (PlantUML starter in the source comments).}}
  \caption{Class diagram of the Document Model.}
  \label{fig:class}
\end{figure}

\section{Use-Case Diagram}
% SOURCES: ch. 5 use cases (UC-1..UC-4)
% GUIDANCE: PlantUML starter (render to figures/usecase-diagram.png):
%   @startuml
%   left to right direction
%   actor Author
%   actor "MCP Client" as MCP
%   rectangle "Draw-to-Web" {
%     Author -- (Author & export a page)
%     Author -- (Draw to create element)
%     Author -- (Match my layout)
%     Author -- (Save / open project)
%     MCP -- (Drive builder headlessly)
%   }
%   @enduml
\begin{figure}[H]
  \centering
  % TODO: \includegraphics[width=0.85\linewidth]{usecase-diagram.png}
  \fbox{\parbox[c][4.5cm][c]{0.9\linewidth}{\centering\itshape
    TODO: use-case diagram for the actors and use cases in ch.~\ref{ch:requirements}.}}
  \caption{Use-case diagram.}
  \label{fig:usecase}
\end{figure}

\section{Sequence Diagrams}
% SOURCES: docs/diagrams/03-export-pipeline.png; src/export/index.ts (9 stages);
%          src/draw/* (draw-to-create flow)
% DRAFT (review): the export pipeline flow is already illustrated.
Figure~\ref{fig:exportflow} shows the nine-stage export pipeline as a flow.

\begin{figure}[H]
  \centering
  \includegraphics[width=0.92\linewidth]{03-export-pipeline.png}
  \caption{The nine-stage export pipeline (validate $\rightarrow$ \ldots $\rightarrow$ save).}
  \label{fig:exportflow}
\end{figure}

% GUIDANCE: add a proper UML SEQUENCE diagram for at least one interaction.
% PlantUML starter for draw-to-create (render to figures/seq-draw.png):
%   @startuml
%   actor Author
%   participant Canvas
%   participant "interpretRectangle" as Interpret
%   participant "snapToGrid" as Snap
%   participant Store
%   Author -> Canvas : drag rectangle
%   Canvas -> Interpret : normalised rect
%   Interpret --> Canvas : ranked kind + confidence
%   Canvas -> Snap : bounds, grid, breakpoint
%   Snap --> Canvas : grid placement
%   Canvas -> Store : insertElement(op)
%   Store --> Canvas : updated document
%   Canvas --> Author : element placed
%   @enduml
\begin{figure}[H]
  \centering
  % TODO: \includegraphics[width=\linewidth]{seq-draw.png}
  \fbox{\parbox[c][4.5cm][c]{0.9\linewidth}{\centering\itshape
    TODO: sequence diagram for the draw-to-create interaction
    (PlantUML starter in the source comments).}}
  \caption{Sequence diagram: draw-to-create.}
  \label{fig:seqdraw}
\end{figure}
